\label{sec:experiments}

\subsection{Experimental Setup}

We compare multi-constraint and edge-weighted optimization across five states with significant minority populations and clear VRA requirements.

\subsubsection{States and Demographics}

Table~\ref{tab:states} shows the five states tested, with district counts, total minority percentages, and target MM district counts based on proportional representation.

\begin{table}[h]
\centering
\caption{States tested with demographics and VRA targets.}
\label{tab:states}
\small
\begin{tabular}{lccc}
\toprule
\textbf{State} & \textbf{Districts} & \textbf{Minority \%} & \textbf{Target MM} \\
\midrule
Alabama & 7 & 36.9\% & 2 \\
Georgia & 14 & 49.9\% & 5 \\
Louisiana & 6 & 44.2\% & 2 \\
Mississippi & 4 & 44.6\% & 2 \\
South Carolina & 7 & 37.9\% & 3 \\
\bottomrule
\end{tabular}
\end{table}

\textbf{Data:} Census 2020 tract-level data including total population and minority voting-age population (VAP). ``Minority'' includes all non-white racial groups per VRA guidelines~\cite{ada1982}.

\subsubsection{Multi-Constraint Configurations}

For each state, we test four minority constraint tolerances:
\begin{itemize}
    \item $\text{ubvec} = [1.005, 1.3]$ (population: $\pm0.5\%$, minority: $\pm30\%$)
    \item $\text{ubvec} = [1.005, 1.5]$ (population: $\pm0.5\%$, minority: $\pm50\%$)
    \item $\text{ubvec} = [1.005, 2.0]$ (population: $\pm0.5\%$, minority: $\pm100\%$)
    \item $\text{ubvec} = [1.005, 5.0]$ (population: $\pm0.5\%$, minority: $\pm400\%$)
\end{itemize}

Target weights specify 60\% minority concentration for target MM districts, with remaining minority population distributed among non-MM districts.

\textbf{Total multi-constraint experiments:} $5 \text{ states} \times 4 \text{ configs} = 20$

\subsubsection{Edge-Weighted Configurations}

Edge-weighted experiments (from prior work~\cite{dellaLibera2026edgeWeighted}) test combinations of:
\begin{itemize}
    \item \textbf{Weight factors ($\alpha$):} 1, 5, 10, 50, 100, 500, 1000
    \item \textbf{Minority thresholds ($\theta$):} 0.40, 0.45, 0.50, 0.55
\end{itemize}

Edges between tracts with $\geq\theta$ minority percentage receive weight $\alpha$; all other edges have weight 1.

\textbf{Total edge-weighted experiments:} $5 \text{ states} \times 7 \text{ factors} \times 4 \text{ thresholds} = 140$

\subsubsection{METIS Parameters}

Both methods use identical METIS settings except for the constraint/weight specifications:
\begin{itemize}
    \item \textbf{Algorithm:} Recursive bisection (\texttt{-ptype=rb})
    \item \textbf{Objective:} Edge cut minimization (\texttt{-objtype=cut})
    \item \textbf{Iterations:} 100 (\texttt{-niter=100})
    \item \textbf{Random seed:} Fixed for reproducibility (\texttt{-seed=42})
\end{itemize}

\subsection{Evaluation Metrics}

\subsubsection{Primary Metrics}

\begin{enumerate}
    \item \textbf{MM District Count:} Number of districts with $\geq50\%$ minority VAP
    \item \textbf{Success:} Configuration succeeds if MM count $\geq$ target
    \item \textbf{Maximum Minority \%:} Highest minority percentage across all districts
    \item \textbf{Edge Cut:} Number of edges crossing district boundaries (unweighted)
\end{enumerate}

\subsubsection{Success Rates}

We compute two success rate metrics:
\begin{itemize}
    \item \textbf{Configuration-level:} Fraction of all configurations achieving target MM count
    \item \textbf{State-level:} Fraction of states where $\geq1$ configuration achieves target
\end{itemize}

\subsection{Reproducibility}

All code, data, and results are available at [ANONYMIZED]. Experiments run on [SYSTEM SPECS: CPU, RAM, OS]. Total runtime: approximately 8 hours for all 160 experiments.
